% Compile with: lualatex ttft_formula_summary.tex
\documentclass[a4paper,11pt]{ltjsarticle}

\usepackage[margin=24mm]{geometry}
\usepackage{amsmath}
\usepackage{booktabs}
\usepackage{hyperref}

\title{Request送信時情報によるTTFT予測式の要約}
\author{LLMServingSim Experiment 2026-07-16}
\date{}

\begin{document}
\maketitle

\section{目的}

Request送信時点で取得可能なprompt、prefix cache、到着負荷、GPUのqueue・KV容量、
routing policyを入力として、end-to-endのTime to First Token（TTFT）を予測する。
TTFTはRouter待ち、Scheduler待ち、Prefill計算、通信時間へ分解する。

\section{TTFTの基本式}

\begin{equation}
\boxed{
\widehat{TTFT}
=
\hat{t}_{\mathrm{route}}
+\hat{t}_{\mathrm{sched}}
+\hat{t}_{\mathrm{compute}}
+\hat{t}_{\mathrm{comm}}
}
\end{equation}

本データでは通信時間が小さく、request送信後のredirect判断で確定するため、point
predictionでは$\hat{t}_{\mathrm{comm}}=0$とする。

\section{各componentの簡易式}

\subsection{Router queue}

Router待ちは、queueが発生する確率と、発生した場合の待ち時間の積で表す。

\begin{align}
p_{\mathrm{route}}
&=
\sigma\!\left(
\beta_0+\boldsymbol{\beta}^{\mathsf T}\boldsymbol{z}
\right), \\
\hat{t}_{\mathrm{route},+}
&=
\exp\!\left[
\operatorname{clip}\!\left(g(\boldsymbol{z}),a,b\right)
\right]-1, \\
\boxed{
\hat{t}_{\mathrm{route}}
&=p_{\mathrm{route}}\hat{t}_{\mathrm{route},+}
}.
\end{align}

ここで$\sigma$はlogistic関数、$\boldsymbol{z}$は標準化済み入力、$g$は深さ2の
回帰木100本からなるpiecewise predictorである。Router queueの発生には、GPUの
capacity pressure、予測KV占有量、空きKV容量、running request数が強く影響する。

\subsection{Scheduler queue}

\begin{equation}
\boxed{
\hat{t}_{\mathrm{sched}}
=
\max\!\left(
0,\gamma_0+\boldsymbol{\gamma}^{\mathsf T}\boldsymbol{z}
\right)
}
\end{equation}

Scheduler待ちには、対象GPUのwaiting request数、直近1秒の到着数、prefix cache量が
主に影響する。

\subsection{Compute / prefill}

Compute時間はinput token数$I$とhome GPU上のcached prefix token数$C$から求める。

\begin{equation}
\boxed{
\hat{t}_{\mathrm{compute}}
=
\max\!\left(
0,
-96.717513
+0.1440215 I
-0.1190117 C
\right)
\quad [\mathrm{ms}]
}
\end{equation}

Input tokenが増えるほどprefill計算量が増え、cached prefixが増えるほど再計算量が減る。

\section{最終的な簡易TTFT式}

以上をまとめると、point predictionは次式になる。

\begin{equation}
\boxed{
\begin{aligned}
\widehat{TTFT}
={}&
\sigma\!\left(
\beta_0+\boldsymbol{\beta}^{\mathsf T}\boldsymbol{z}
\right)
\left\{
\exp\!\left[
\operatorname{clip}\!\left(g(\boldsymbol{z}),a,b\right)
\right]-1
\right\} \\
&+
\max\!\left(
0,\gamma_0+\boldsymbol{\gamma}^{\mathsf T}\boldsymbol{z}
\right) \\
&+
\max\!\left(
0,-96.717513+0.1440215I-0.1190117C
\right).
\end{aligned}
}
\end{equation}

実装で使用する定数は、

\begin{equation}
\beta_0=-14.013232,
\qquad
\gamma_0=60.427288\ \mathrm{ms},
\qquad
(a,b)=(0.001642,10.878924)
\end{equation}

である。$\boldsymbol{\beta}$、$\boldsymbol{\gamma}$、標準化用の平均・scale、回帰木の
全係数はCSVおよびJSON artifactに保存されている。

\section{主要入力の意味}

\begin{center}
\begin{tabular}{lp{88mm}}
\toprule
入力群 & 意味 \\
\midrule
Request size & Input/output token数とcached prefix token数 \\
Arrival load & 平均request rate、直近1秒・5秒のglobal/home到着数 \\
Queue state & 対象GPUのwaiting request数とrunning request数 \\
KV capacity & Required KV、予測active KV、available KV、capacity pressure \\
Candidate state & 収容可能GPU数、全候補のqueue・KV状態の合計または最小・最大 \\
Policy & Local待機、cold redirect、KV handoffのone-hot表現 \\
\bottomrule
\end{tabular}
\end{center}

特にRouter queueの中心的な要因は、送信時点でrequestをKV容量とsequence slotの両方に
ついて収容できるGPUが存在するかどうかである。

\section{評価結果と制約}

13,500 requests、15 independent scenariosのleave-one-scenario-out評価では、TTFTの
MAEは716.1\,ms、$R^2$は0.467であった。通常領域のcomputeとqueue発生判定は比較的
良好だが、高負荷時の長いRouter waitを過小予測する。

Routingの安全判定では、正値Router待ちへscenario-held-out residualの90 percentileを
加えた上側予測も用意されている。

\begin{equation}
\hat{t}_{\mathrm{route,upper}}
=
\hat{t}_{\mathrm{route},+}
+10{,}311.69\ \mathrm{ms}
\end{equation}

この固定補正は長時間待機の見逃しを防ぐ一方、1秒程度のlocal待機上限と組み合わせると
redirect判断をほぼ常時ONにする場合がある。したがって、point prediction、SLA上側予測、
候補GPUランキングは目的別に分けて利用する必要がある。

\end{document}
